\chapter{Other applications of tensorization techniques}%
\label{parOtherApplicationsOfTensorizationTechniques}%
\addcontentsline{itc}{chapter}{\protect\numberline{\thechapter}Other applications of tensorization techniques}

In the previous chapter we have been seeing
how Hilbertian decorrelation hypotheses between pairs of variables
could yield `global' results on an arbitrary number of variables,
by splitting functions of several variables into relevant telescopic sums.
I used the word ``tensorization'' to qualify these results,
as the conclusions were of the same nature as the hypotheses.

But the techniques of \S~\ref{parTensorization} can also be applied to get other types of results.
In this chapter I am going to show how, from Hilbertian decorrelation hypotheses,
one can get results on some classical features of particle systems
which are not linked with Hilbertian correlations \emph{a priori}.

I will deal with two such features.
First, I will look at the implications of $\rho$-mixing
on the existence of a central limit theorem%
—more precisely, of a \emph{spatial} central limit theorem,
since I am more interested in random fields than in sequences
(variables indexed by~$\ZZ^n$ rather than by~$\ZZ$).
Very sharp results concerning this issue are already known;
however, I find interesting to show how it goes with my `tensorization-like' approach:
this approach takes indeed a quite different way to do the job,
which may be neater by certain sides.
Moreover, the results are stated with a slighlty different vocabulary%
—namely, \emph{subjective} Hilbertian correlations.

Next, I will look at the question of spectral gap for Glauber dynamics.
Though this point has already been thouroughly studied in a $\beta$-mixing paradigm,
this work, to the best of my knowledge,
is the first to show how $\rho$-mixing can be used to tackle this issue.

My main goal here is just to show
\emph{how} the techniques of this work may be applied to the problems
of spatial central limit theorem and convergence of the Glauber dynamics.
Accordingly, I favoured the simplicity on proofs against the refinement of the results.



\section{Spatial central limit theorem}

\subsection{Introduction}

A fundamental result in probability theory is the central limit theorem (CLT),
which, in its standard statement, requires an assumption of complete independence.
It is natural to wonder whether that assumption can be relaxed
into an hypothesis of `near independence'.
Hilbertian decorrelations are a natural frame for such a generalization,
since the CLT already takes place in an $L^2$ setting.

Our point of view is motivated by statistical physics.
Let~$\ZZ^n$ be a lattice, on each vertex $i$ of which there is a random `spin' $X_i$
ranged in some space $\mcal{X}$ not depending on~$i$.
We assume that the law of the system is translation invariant,
i.e.\ that for all~$z\in\ZZ^n$, $(X_{i+z})_{i\in\ZZ^n}$ has the same law
as~$\vec{X}_{\ZZ^n}$.
Then, for all~$z \in \ZZ^n$, we denote
\[\label{f5933} \epsilon_z = \{ X_i : X_{i+z} \}_* .\]
We are interested in situations where the~$\epsilon_z$ are sufficiently `rapidly decreasing'
as $|z| \longto \infty$ so that $\sum_{z\in\ZZ^n} \epsilon_z < \infty$.

Let~$f \colon \mcal{X} \to \RR$ be a function such that $f(X_0)$ is square-integrable and centered.
The question is, does one get a CLT when summing $f(X_i)$ for~$i$ in a large subset of~$\ZZ^n$,
i.e., does the sum grow as the square root of the number of its terms
and have asymptotically normal distribution?
For instance, we would like the law of the variable
\[ \frac{1}{\sqrt{\ell^n}} \sum_{\substack{i\in\ZZ^n\\0\leq i_1,\ldots,i_n<\ell}} f(X_i) \]
to weakly converge, when~$\ell \longto \infty$, to some Gaussian distribution.

\begin{Rmk}
Note that the limit distribution, if it exists, will have to be centered,
but its variance will not be equal to~$\Var\big(f(X_0)\big)$ in general.
\end{Rmk}

In the case $n=1$, extremely sharp results for this topic have been known from long;
let us cite, among many others, \cite{Rosenblatt56, Ibragimov75, Peligrad87, Bradley92}.
For $n\geq 2$, similar results also exist; see e.g.\ \cite[Theorem~5]{Bradley92} for such a result,
and~\cite[\S~29]{Bradley-book} for a survey of the topic.
All these proofs relie on some `coupling' between (bunches of) the spins
and other convenient variables which are close to them,
but which are \emph{actually} independent,
so as to deduce the CLT for the former from the CLT for the latter.
On the other hand, my proof will mimick
Lévy's proof of the CLT, hence needing no coupling argument.

\emph{A priori} the results presented here do not improve the state of the art;
however, when turning to quantitative versions of these results,
it is likely that the difference between the usual method and mine
would yield a difference in the corresponding non-asymptotic bounds obtained.


\subsection{Product of weakly coupled variables}

My results relie on the following
%
\begin{Lem}\label{l1840}
Let~$N \geq 1$ and let~$\dot{\mcal{F}}_1, \ldots, \dot{\mcal{F}}_N$ be $\sigma$-algebras with
$\{ \dot{\mcal{F}}_i : \dot{\mcal{F}}_j \}_* \leq \epsilon_{ij}$%
\footnote{As in \S~\ref{parMachineryForUsingTheTensorizationTheorems},
``$*$'' stands for ``the natural $\sigma$-metalgebra of the system''.
Moreover, in the same way as in \S~\ref{parMinimalHypotheses},
it is actually possible in the statement of the lemma
to replace that~$\sigma$-metalgebra by smaller ones.},
and denote
\[ \bar\epsilon = \sup_i \sum_{j\neq i} \epsilon_{ij} .\]
Let~$\Phi_1, \ldots, \Phi_N$ be complex-valued random variables with $|\Phi_i|\leq1$ a.s.,
such that $\Phi_i$ is $\dot{\mcal{F}}_i$-measurable for all~$i$,
with all the~$\Phi_i$ having the same distribution.
Then, denoting by~$\phi$ the common value of the~$\EE[\Phi_i]$,
\[\label{fo9356} \Big| \EE\Big[ \prod_{i} \Phi_i \Big] - \phi^N \Big| \leq
N \bar\epsilon(1+\bar\epsilon) (1-|\phi|^2) .\]
\end{Lem}

\begin{proof}
Denote $\delta \coloneqq \ecty(\Phi)$. Since $\EE[|\Phi^2|] \leq 1$,
the definition of (complex) variance ensures that $\delta \leq \sqrt{1-|\phi|^2}$.

For all~$i \in \{0,\ldots,N\}$, denote $\mcal{F}_i \coloneqq \bigvee_{i'\leq i}\dot{\mcal{F}}_i$ ;
denote $\Psi^{(i)} \coloneqq \prod_{i'\leq i} \Phi_{i'}$; define
\[ \Psi^{(i)}_j \coloneqq (\Psi^{(i)})^{\mcal{F}_j} - \EE[\Psi^{(i)}|\mcal{F}_{j-1}] \]
and denote $\Delta^{(i)}_j \coloneqq \ecty(\Psi^{(i)}_j)$.
Also denote
\[ \Phi_{i,j} \coloneqq (\Phi_i)^{\mcal{F}_j} - \EE[\Phi_i|\mcal{F}_{j-1}] .\]

Usual manipulation on conditioning shows that, for~$i\geq1$,
\[ \Psi^{(i)}_j = \Psi^{(i-1)}_j (\Phi_i)^{\mcal{F}_j}
+ \Psi^{(i-1)}_{j-1} \Phi_{i,j} .\]
Since $\|\Phi_i\|_{L^\infty} \leq 1$,
one has also $\| (\Phi_i)^{\mcal{F}_j} \|_{L^\infty} \leq 1$, hence
\[ \ecty\big( \Psi^{(i-1)}_j (\Phi_i)^{\mcal{F}_j} \big)
\leq \ecty(\Psi^{(i-1)}_j) = \Delta^{(i-1)}_j .\]
%
Similarly, it is obvious that $\|\Psi^{(i-1)}\|_{L^\infty} \leq 1$, whence
\[ \ecty\big( \Psi^{(i-1)}_{j-1} \Phi_{i,j} \big) \leq \ecty(\Phi_{i,j}) .\]

Now, I claim that
%
\begin{Clm}\label{clm-dld}
\[ \ecty(\Phi_{i,j}) \leq \epsilon_{ij} \delta .\]
\end{Clm}
%
\begingroup\def\proofname{Proof of Claim~\ref{clm-dld}}
\begin{proof}
Conditionally to~$\mcal{F}_{j-1}$, $\Phi_{i,j}$ is indeed the projection on~$\dot{\mcal{F}}_{j}$ of the centered $\dot{\mcal{F}}_i$-measurable function $(\Phi_i - \EE[\Phi_i|\mcal{F}_{j-1}])$, whose standard deviation is less than $\ecty(\Phi_i) = \delta$ by associativity of the variance, so that $\ecty(\Phi_{i,j}) \leq {\{\dot{\mcal{F}}_i:\dot{\mcal{F}}_j\}_{\mcal{F}_{j-1}}} \delta \ab \leq \epsilon_{ij} \delta$.
\end{proof}\endgroup

In the end, we got that
\[ \Delta^i_j \leq \Delta^{(i-1)}_j + \epsilon_{ij} \delta .\]
Since $\Delta^0_j = 0$, one has therefore:
\[ \forall i,j \qquad \Delta^i_j \leq (1+\bar\epsilon) \delta .\]
%
Now, denoting $\psi^{(i)} \coloneqq \EE[\Psi^{(i)}]$, one has
\[ | \psi^{(i)} - \phi\psi^{(i-1)} |
= \Big| \sum_{j<i} \EE\big[ \Psi^{(i-1)}_j \Phi_{i,j} \big] \Big|
\leq \sum_{j<i} \Delta^{(i-1)}_j \ecty(\Phi_{i,j})
\leq \bar\epsilon (1+\bar{\epsilon}) \delta^2 ,\]
%
and finally
\[ | \psi^{(N)} - \phi^N |
= \sum_{i=1}^N |\phi|^{N-i} \big|\psi^{(i)} - \phi\psi^{(i-1)}\big|
\leq \sum_{i=1}^N \big|\psi^{(i)} - \phi\psi^{(i-1)}\big| \\
\leq N \bar\epsilon(1+\bar\epsilon) \delta^2 ,\]
which is (\ref{fo9356}) if you recall that $\delta^2 \leq 1-|\phi|^2$.
\end{proof}


\subsection{A spatial CLT}

First I state and prove a CLT on cubes:
%
\begin{Thm}\label{thm0318}
Consider a translation-invariant spin model on a lattice $\ZZ^n$
and define $\epsilon_z$ by~(\ref{f5933}).
Assume that $\sum_{z\in\ZZ^n} \epsilon_z < \infty$.
Then for any centered square-summable function~$f \colon \mcal{X} \to \RR$,
there exists a constant $\sigma < \infty$ such that
\[\label{f7877}
\frac{1}{\sqrt{\ell^n}} \sum_{\substack{i\in\ZZ^n\\0\leq i_1,\ldots,i_n<\ell}} f(X_i)
\stackrel{\ell\longto\infty}{\rightharpoonup} \mcal{N}(\sigma^2) ,\]
where ``$\rightharpoonup$'' denotes convergence in law.
\end{Thm}

\begin{proof}
Denote by~$F(\ell)$ —or merely $F$ —the left-hand side of~(\ref{f7877}).

What will be the value of~$\sigma$?
Clearly we must have
\[\label{fo0116} \sigma^2 = \lim_{\ell\longto\infty} \Var \big( F(\ell) \big) ,\]
which yields
\[ \sigma = \sqrt{ \sum_{z\in\ZZ} \EE\big[ f(X_0) f(X_z) \big] } ,\]
where the expression under the root sign,
which is necessarily nonnegative,
is finite because $|\EE[f(X_0)f(X_z)]| \leq
\epsilon_z \ecty(f(X_0))\ecty(f(X_z)) = \epsilon_z \Var(f)$.
By the way, we will denote
\[ \sigma_*^2 \coloneqq \Big( \sum_z \epsilon_z \Big) \|f\|_{L^2}^2 .\]

Fix some arbitrary $\eta > 0$.
The assumption that $\sum \epsilon_z < \infty$
implies the existence of an $\ell_0 < \infty$
such that
\[ \sum_{|z|_\infty > \ell_0} \epsilon_z \leq \eta ,\]
where $|z|_\infty$ denotes $\max(|z_1|,\ldots,|z_n|)$.
By~(\ref{fo0116}), we can also fix an $\ell_1 < \infty$
such that
\[\label{f1939} \big| \Var \big( F(\ell) \big) - \sigma^2 \big| \leq \eta .\]

Now we will `tile' the cube of size $\ell$ into a `patchwork'
made of cubes of size $\ell_1$ which I call ``tiles'',
each tile being at distance at least $\ell_0$ from the others, plus some ``scrap''.
I denote by~$\tilde{F}$ the part of~$F$ due to the tiles
and by~$F^*$ the part of~$F$ due to the scrap.

Index the tiles by~$\{ 1,\ldots,N \}$, with
$N \coloneqq \big\lfloor (\ell+\ell_0) \div (\ell_1+\ell_0) \big\rfloor^n$.
We write, with obvious notation, $\tilde{F} \eqqcolon F_1 + \cdots + F_N$.
For~$\lambda\in\RR$, denote
\begin{eqnarray}
\Psi (\lambda,\ell) \coloneqq \exp(i\lambda \tilde{F}), &\qquad&
\psi(\lambda,\ell) \coloneqq \EE[\Psi (\lambda,\ell)] ; \\
\Phi_j(\lambda,\ell) \coloneqq \exp(i\lambda F_j), &\qquad&
\phi(\lambda,\ell) \coloneqq \EE[\Phi_j(\lambda,\ell)] .
\end{eqnarray}
%
Then we are exactly in situation of applying Lemma~\ref{l1840},
which yields:
\[\label{f1890} \big| \psi(\lambda,\ell) - \phi(\lambda,\ell)^N \big| \leq
N \eta(1+\eta) \big( 1-|\phi(\lambda,\ell)|^2 \big) .\]

Let us look at the asymptotics of Formula~(\ref{f1890}) when~$\ell \longto \infty$.
We observe that, denoting
\[ F_{\rm t} \coloneqq \frac{1}{\sqrt{\ell_1^n}} \sum_{i\in\text{\it\ fixed tile}} f(X_i) ,\]
one has
\[ F_j \stackrel{\text{law}}{=} \frac{\sqrt{\ell_1^n}}{\sqrt{\ell^n}} F_{\rm t} .\]
%
Since $F_{\rm t}$ is centered, its Fourier transform satisfies $\hat{F}_{\rm t}(0)=1$,
$\hat{F}_{\rm t}'(0)=0$ and $\hat{F}_{\rm t}'' = \Var(F_{\rm t})$, so that
\[ \lim_{\ell\longto\infty} \ell^n \big(1-\phi(\lambda,\ell)\big)
= \frac{\lambda^2}{2} \ell_1^n \Var(F_{\rm t}) ,\]
where, denoting $\sigma^2_{\ell_1} \coloneqq \Var(F_{\rm t})$,
we recall that $\ell_1$ has been taken sufficiently large so that
$|\sigma^2_{\ell_1} - \sigma^2| \leq \eta$.
%
Then, since $N \stackrel{\ell\longto\infty}{\sim} \ell^n \div (\ell_1+\ell_0)^n$,
one has the following asymptotics for~(\ref{f1890}):
\begin{eqnarray}
\phi(\lambda,\ell)^N & \stackrel{\ell\longto\infty}{\longto} &
\exp \left[ -\sigma^2_{\ell_1} \bigg( \frac{\ell_1}{\ell_1+\ell_0} \bigg)^{\!n}
\frac{\lambda^2}{2} \right] ; \\
N \big( 1-|\phi(\lambda,\ell)|^2 \big) & \stackrel{\ell\longto\infty}{\longto} &
\sigma^2_{\ell_1} \bigg( \frac{\ell_1}{\ell_1+\ell_0} \bigg)^{\!n}\lambda^2 .
\end{eqnarray}

It remains to control the contribution of~$F^*$.
%
\begin{Clm}\label{c7901}
There are at most
\[ \bigg[1-\bigg(\frac{\ell_1}{\ell_1+\ell_0}\bigg)^{\!n}\bigg]\ell^n + n\ell_1\ell^{n-1} \]
scrap spins.
\end{Clm}
%
By Claim~\ref{c7901},
\[ \|F^*\|_{L^1} \leq \Var(F^*) \leq
\bigg[1-\bigg(\frac{\ell_1}{\ell_1+\ell_0}\bigg)^{\!n} + n\frac{\ell_1}{\ell} \bigg] \sigma_*^2 ;\]
then the contribution of~$F^*$ is controlled using the following immediate
%
\begin{Lem}
Let~$X$ and~$H$ be real random variables with $\|H\|_{L^1} < \infty$.
Then, for~$\lambda\in\RR$,
\[ \big| \EE\big[e^{i\lambda(X+H)}\big] - \EE\big[e^{i\lambda X}\big] \big| \leq |\lambda| \|H\|_{L^1} .\]
\end{Lem}

In the end, putting everything together we get:
\begin{multline}\label{f6829}
\limsup_{\ell\longto\infty} \big| \psi(\lambda,\ell) - e^{-\sigma^2\lambda^2/2} \big| \leq \\
\Big|e^{\big{-(\sigma^2-\eta) [\ell_1 \div (\ell_1+\ell_0)]^n \lambda^2/2}} - e^{-\sigma^2\lambda^2/2}\Big|
+ \eta(1+\eta) \bigg( \frac{\ell_1}{\ell_1+\ell_0} \bigg)^{\!n} (\sigma^2+\eta) \lambda^2
+ \sqrt{ 1-\bigg(\frac{\ell_1}{\ell_1+\ell_0}\bigg)^{\!n} } \sigma_* .
\end{multline}
%
Since there were no upper restriction on the value of~$\ell_1$, we can assume
that we have taken it such that $[\ell_1 \div (\ell_1+\ell_0)]^n \geq 1-\eta$.
Then (\ref{f6829}) becomes:
\[\label{f6890} \limsup_{\ell\longto\infty}
\big| \psi(\lambda,\ell) - e^{-\sigma^2\lambda^2/2} \big| \leq
\Big| e^{\big{-(\sigma^2-\eta)(1-\eta)\lambda^2/2}} - e^{-\sigma^2\lambda^2/2} \Big|
+ \eta(1+\eta)(1-\eta) (\sigma^2+\eta) \lambda^2
+ \sqrt{\eta} \sigma_* .\]
%
The right-hand side of~(\ref{f6890}) can be made arbitrarily close to~$0$
by taking $\eta$ small enough, so we have proved that
%
\[ \forall \lambda \in \RR \quad \EE\big[e^{i\lambda F(\ell)}\big]
\stackrel{\ell\longto\infty}{\longto} e^{-\sigma^2\lambda^2/2} .\]
By L\'evy's theorem on characteristic functions,
this is tantamount to saying that $F(\ell)$ converges in law to~$\mcal{N}(\sigma^2)$.
\end{proof}

The CLT should remain valid for other shapes than a cube,
since morally the random field $f(X_i)$
should look like a Gaussian white noise at large scales.
Indeed, the same proof as above yields a CLT for general shapes,
where moreover convergence is uniform in the shape considered in some way.
Let us give a precise statement:
%
\begin{Def}
An open subset $U \subset \RR^n$ (not necessarily connected)
is said to be \emph{$\mcal{C}^2$} if its boundary $M$ is a $\mcal{C}^2$ submanifold
of~$\RR^n$ (of codimension $1$).
We define the \emph{roughness} of~$U$, denoted by~$\kappa(U)$, as
\[ \kappa(U) \coloneqq \sup_{x\in M} \VERT \II(x) \VERT, \]
where $\II(\Bcdot)$ denotes the shape tensor of~$M$~\cite[Chapter~10]{Differential},
which measures the local deviation of~$M$ from being flat.
Also, the Lebesgue measure of~$U$ will be denoted by~$\volume(U)$.
\end{Def}

\begin{Thm}\label{thm3015}
Consider a translation-invariant spin model on a lattice $\ZZ^n$
and define $\epsilon_z$ by~(\ref{f5933}).
Assume that $\sum_{z\in\ZZ^n} \epsilon_z < \infty$.
Then for any centered square-summable function~$f \colon \mcal{X} \to \RR$,
if $(U_k)_{k\in\NN}$ is a sequence of $\mcal{C}^2$ bounded subsets of~$\RR^n$
with $\sup_k \kappa(U_k) < \infty$ and
$(\ell_k)_{k\in\NN}$ is a sequence of positive numbers tending to infinity,
\[ \frac{1}{\sqrt{\ell_k^n\volume(U_k)}} \sum_{i\in\ell_kU_k\cap\ZZ^n} f(X_i)
\stackrel{k\longto\infty}{\rightharpoonup} \mcal{N}(\sigma^2) ,\]
where $\sigma^2$ is the same as in Theorem~\ref{thm0318}.
\end{Thm}

\begin{proof}
Just copy the proof of Theorem~\ref{thm0318}.
The only difference lies in proving the analoguous of Claim~\ref{c7901},
which is where one needs the~$\kappa(U_k)$ to be bounded.
Observe that we use the non-asymptotic form of our intermediate bounds
to get a result independent of the precise shape of the~$U_k$.
\end{proof}

\begin{Rmk}
Another generalization of the CLT, still based on the idea that
the field $f(X_i)$ looks like a Gaussian white noise at large scales,
is the statement that for~$\phi$ a continuous function with compact support,
\[ \frac{1}{\sqrt{\ell^n}} \sum_{i\in\ZZ^n} \phi\big(X_i\div\ell\big) f(X_i)
\stackrel{\ell\longto\infty}{\rightharpoonup}
\mcal{N}\big(\sigma^2 \mathsmaller{\int_{\RR^n} \phi(x)^2\,\dx{x}} \big) .\]
This can be proved with the same methods as before.
\end{Rmk}


\section{Spectral gap for the Glauber dynamics}

\subsection{Introduction}

In this section we are looking at a probabilistic system
made of a large number of `elementary' random variables $(X_i)_{i\in I}$
— $I$ may be seen as lattice and $X_i$ as the state of the particle being at site~$i$.
As is customary by now, theorems will only be stated in the case where $I$ is finite,
the infinite case being got by passing to the limit.

\begin{Def}
Denoting by~$\Omega$ the states space of~$\vec{X}_I$,
let~$\Pr$ be a probability measure on~$\Omega$.
The \emph{Glauber dynamics}~\cite{Glauber, Dobrushin71}
associated to~$\Pr$ is the Markov process on~$\Omega$ having the following law:
on each~$i\in I$ there is an alarm clock,
all the clocks being independent and ringing with law $\mathit{Poisson}(1)$.
When a clock rings, the state of spin $X_i$ —and only it—is flipped
so that the state of~$X_i$ immediately after the flip follows the law
$\Pr\big(X_i|\vec{X}_{I\setminus\{i\}}\big)$.

In formal terms, the Glauber dynamics is the Markov process
whose generator $\mcal{L}$ on~$L^\infty(\Omega)$ is defined by:
\[ \big(\mcal{L}f\big)(\vec{x}_I)
= \sum_{i\in I} \EE \big[ f(\vec{X}_I) - f(\vec{x}_I) \big|
\vec{X}_{I\setminus\{i\}} = \vec{x}_{I\setminus\{i\}} \big] .\]
\end{Def}

Let us recall some basic facts on the Glauber dynamics
(see~\cite[Chapter~IV]{Liggett} for more details).
By construction $\Pr$ is a reversible equilibrium measure for the dynamics,
so $\mcal{L}$ is self-adjoint on~$L^2(\Pr)$.
Since obviously $\mcal{L}1\equiv0$, one can also define $\mcal{L}$ on~$\ldb(\Pr)$,
on which it is self-adjoint too. This leads to the following definition:
%
\begin{Def}
The \emph{energy} of~$f\in\ldb(\Pr)$ is
\[ \mcal{E}(f,f) = \langle Lf, f \rangle \]
\end{Def}
%
The following immediate identity shows that $\mcal{E}$ is always a nonnegative bilinear form:
\begin{Pro}\label{p1345}
\[ \mcal{E}(f,f) = \int_\Omega \dx\Pr[\vec{x}_I]
\sum_i \Var\big( f \big| \vec{X}_{I\setminus\{i\}} = \vec{x}_{I\setminus\{i\}} \big) .\]
\end{Pro}

\begin{Def}
For~$\lambda > 0$, the Glauber dynamics is said to have \emph{spectral gap $\geq\lambda$}
if, for all~$f\in\ldb(\Pr)$,
\[\label{fo8103} \mcal{E}(f,f) \geq \lambda \Var(f) .\]
\end{Def}
%
What makes spectral gap interesting is that its positiveness
is equivalent to exponential convergence to~$0$
of the semigroup $(e^{-t\mcal{L}})_{t\geq0}$ on $\ldb(\Pr)$,
the rate of convergence being equal to the width of the spectral gap.
As the Glauber dynamics is one of the easiest ways to simulate the law $\Pr$ for complicated models,
the stake of having exponential convergence for it is evident.

Many works have been done on the spectral gap of the Glauber dynamics,
see for instance Martinelli's St-Flour course~\cite{Martinelli}.
Several results state that,
the less spins are correlated, the larger the spectral gap is.
Yet the researchers who work on this topic generally express
the decorrelation between the spins in terms of $\beta$-mixing (cf.\ Definition~\ref{def3167}),
while it seems be more natural to look at them in terms of Hilbertian decorrelations,
since the formula~(\ref{fo8103}) stating the spectral gap problem
takes place in a Hilbertian frame itself.
Thus my goal here will be to find a control on the spectral gap
expressed in terms of $\rho$-mixing conditions.
Since Hilbertian correlations look to be the minimal frame
to study the spectral gap for the Glauber dynamics,
hopefully the bounds yielded by this method will be sharp.

Another noticeable feature of my approach
is that it remains at a quite abstract level:
no symmetry property of~$I$ or $\Pr$ need be assumed,
all the work essentially consisting in manipulating relevant quadratic forms.



\subsection{A lower bound for the spectral gap}

The central theorem of this section is the following:
\begin{Thm}\label{t7230}
Take $I = \{1,\ldots,N\}$.
Suppose that for all distinct $i,j\in I$ one has ${\{X_i:X_j\}_\ast} \ab \leq \epsilon_{ij} < 1$
—we will make the costless assumption that $\epsilon_{ji} = \epsilon_{ij}$.
For~$i \in I$, denote
\[ \tilde1_i \coloneqq \frac{1}{ \prod_{i<j\leq N} (1-\epsilon_{ij}^2) }
= \frac{\teps_{iN}}{\epsilon_{iN}} ,\]
and for~$i < j$, denote
\[ \teps_{ij} \coloneqq
\frac{\epsilon_{ij}}{\prod_{i<j'\leq j}(1-\epsilon_{ij'}^2)} .\]
Then the Glauber dynamics has spectral gap at least $\VERT M \VERT^{-2}$,
where $M$ is the $(N \times N)$ matrix defined by
\[\label{for4393} M = \begin{pmatrix} 1 & -\teps_{12} & \cdots & -\teps_{1N} \\
0 & \ddots & \ddots & \vdots \\
\vdots & \ddots & \ddots & -\teps_{(N-1)N} \\
0 & \cdots & 0 & 1 \end{pmatrix}^{-1}
\begin{pmatrix} \tilde1_1 & 0 & \cdots & 0 \\
0 & \ddots & \ddots & \vdots \\
\vdots & \ddots & \ddots & 0 \\
0 & \cdots & 0 & \tilde1_N \end{pmatrix} .\]
\end{Thm}
%
\begin{Rmk}\label{rmk4123}
The form of the first matrix in the right-hand side of~(\ref{for4393})
ensures that it is invertible.
Since moreover all the~$\epsilon_{ij}$ were supposed $<1$,
all the~$\tilde\epsilon_{ij}$ and the~$\tilde{1}_i$ are finite;
thus, the lower bound~$\VERT M \VERT^{-2}$ is strictly positive.
\end{Rmk}

\begin{proof}
Let~$f$ be a centered square-integrable function on~$(\Omega,\Pr)$.
For~$I' \subset I$, denote $\mcal{F}_{I'} \coloneqq \sigma\big(\vec{X}_{I'}\big)$.
For~$i\in I$, $I' \subset I\setminus\{i\}$,
denote
\[ f_i^{I'} \coloneqq f^{\mcal{F}_{I'\uplus\{i\}}} - \EE[f|\mcal{F}_{I'}] ;\]
define moreover
\begin{eqnarray}
f_i^{\neq} & \coloneqq & f_i^{I\setminus\{i\}} ; \\
f_i^{<}  & \coloneqq & f_i^{\{1,\ldots,i-1\}}.
\end{eqnarray}
%
Then by Proposition~\ref{p1345}, one has
\[ \mcal{E}(f,f) = \sum_i \Var(f_i^{\neq}) ,\]
while the usual telescopic argument shows that
\[ \Var(f) = \sum_i \Var(f_i^<) .\]

So to prove the theorem, we have to establish links between
the different values $\Var(f_i^{I'})$.
It will be convenient to introduce the shorthands $\Delta_i^{I'} = \ecty(f_i^{I'})$.
One has the following
%
\begin{Clm}\label{l8326}
For~$I' \subset I$ and~$i,j\in I \setminus I'$ with $j\neq i$,
\[\label{f4464} \Delta_i^{I'} \leq \Delta_i^{I'\uplus\{j\}} + \epsilon_{ij} \Delta_j^{I'} .\]
\end{Clm}

\begin{proof}
Assume in a first time that $I'=\emptyset$,
and denote $f_i \coloneqq f_i^{\emptyset}$, $f_j \coloneqq f_j^{\emptyset}$, $f_i^j \coloneqq f_i^{\{j\}}$
and $\mcal{F}_i \coloneqq \mcal{F}_{\{i\}}$.
%
Projecting the decomposition ``$f_i = f_i^j + (f_i-f_i^j)$''
on~$L^2(\mcal{F}_i)$, one has $f_i = (f_i^j)^{\mcal{F}_i} + (f_j)^{\mcal{F}_i}$,
whence by the Cauchy--Shwarz inequality:
\[\label{f4448}
\ecty(f_i) \leq \ecty\big((f^j_i)^{\mcal{F}_i}\big) + \ecty\big((f_j)^{\mcal{F}_i}\big) .\]
%
One has trivially
$\ecty\big((f^j_i)^{\mcal{F}_i}\big) \leq \ecty(f^j_i)$;
on the other hand, $f_j$ is $X_j$-measurable,
so $\ecty\big((f_j)^{\mcal{F}_i}\big) \leq \epsilon_{ij} \ecty(f_j)$.
In the end, (\ref{f4448}) becomes
\[ \ecty(f_i) \leq \ecty(f^j_i) + \epsilon_{ij} \ecty(f_j) ,\]
which is~(\ref{f4464}) for $I' = \emptyset$.

In the case $I' \neq \emptyset$, the same reasoning can be performed, except that
one have to work conditionally to~$\mcal{F}_{I'}$.
Then, taking
$f_i = f^{I'}_i$,
$f_j = f^{I'}_j$,
$f^j_i = f^{I'\uplus\{j\}}_i$,
$\mcal{F}_i = \mcal{F}_{I'\uplus\{i\}}$,
%
one gets
\[\label{fo5917}
\ecty(f_i|\mcal{F}_{I'}) \leq \ecty(f^j_i|\mcal{F}_{I'}) + \epsilon_{ij} \ecty(f_j|\mcal{F}_{I'}) .\]
Now
\[\ecty(f_i) = \sqrt{\int \ecty\big(f_i|\vec{X}_{I'}=\vec{x}_{I'}\big)^2 \dx\Pr[\vec{x}_{I'}]} ,\]
with similar formulas for~$f_j$ and~$f^j_i$, since all these functions
are centered w.r.t.~$\mcal{F}_{I'}$.
%
Therefore, integrating~(\ref{fo5917}) and applying Minkowski's inequality yields:
\[ \ecty(f_i) \leq \ecty(f^j_i) + \epsilon_{ij} \ecty(f_j) ,\]
i.e.~(\ref{f4464}).
\end{proof}

For~$i \leq j$, let us denote
\[ \Delta^{[j]}_i = \Delta^{\{1,\ldots,j\} \setminus \{i\}}_i .\]
Claim~\ref{l8326} will be used through the following corollary:
\begin{Clm}\label{l3526}
For all~$i<j$,
\[\label{f7631} \Delta^{[j-1]}_i \leq \frac{1}{1-\epsilon_{ij}^2}
\big( \Delta^{[j]}_i + \epsilon_{ij} \Delta^{<}_j \big) .\]
\end{Clm}

\begin{proof}
We have to bound $\Delta^{[j-1]}_i$, which here we rather denote
$\Delta^{[b-1]}_a$ to avoid confusion with the notation of Claim~\ref{l8326}.
Applying Claim~\ref{l8326} with $I' = \{1,\ldots,b-1\} \setminus \{a\}$,
$i = a$ and $j = b$, one has
\[\label{f7625a} \Delta^{[b-1]}_a = \Delta^{\{1,\ldots,b-1\}\setminus\{a\}}_{a}
\leq \Delta^{\{1,\ldots,b\}\setminus\{a\}}_{a}
+ \epsilon_{ab} \Delta^{\{1,\ldots,b-1\}\setminus\{a\}}_b
= \Delta^{[b]}_{a} + \epsilon_{ab} \Delta^{\{1,\ldots,b-1\}\setminus\{a\}}_b .\]
But applying again Claim~\ref{l8326}, this time with $I' = \{1,\ldots,b-1\} \setminus \{a\}$,
$i = b$ and $j = a$, one has
\[\label{f7625b} \Delta^{\{1,\ldots,b-1\}\setminus\{a\}}_{b}
\leq \Delta^{\{1,\ldots,b-1\}}_{b}
+ \epsilon_{ab} \Delta^{\{1,\ldots,b-1\}\setminus\{a\}}_a
= \Delta^{<}_b + \epsilon_{ab} \Delta^{[b-1]}_a .\]
Combining~(\ref{f7625a}) and~(\ref{f7625b}) then yields~(\ref{f7631}).
\end{proof}

Now let us show how Claim~\ref{l3526} implies the theorem.
To avoid heavy formalism, I will detail the computations for $I = \{1,2,3,4\}$
(rather denoted by~$I = \{a,b,c,d\}$ here to avoid confusions
with ``$1$'' and~``$2$'' taken as numbers),
hoping that generalizing is obvious then.

First, note that
\[\label{f4236d} \Delta^<_d = \Delta^{\neq}_d .\]
Now, by a direct use of Claim~\ref{l3526},
\[\label{f4236c} \Delta^<_c = \Delta^{[c]}_c
\leq \frac{1}{1-\epsilon_{cd}^2} \big( \Delta^{[d]}_c + \epsilon_{cd} \Delta^{<}_d \big)
= \tilde1_c \Delta^{\neq}_c + \tilde{\epsilon}_{cd} \Delta^{\neq}_d .\]
%
To bound $\Delta^<_b$, we have to iterate Claim~\ref{l3526} twice:
\begin{multline}\label{f4236b}
\Delta^<_b = \Delta^{[b]}_b
\leq \frac{1}{1-\epsilon_{bc}^2} \Delta^{[c]}_b + \teps_{bc} \Delta^{<}_c \\
\leq \frac{1}{(1-\epsilon_{bc}^2)(1-\epsilon_{bd}^2)}
\big( \Delta^{[d]}_b + \epsilon_{cd} \Delta^{<}_d \big) + \teps_{bc} \Delta^<_c
= \tilde1_b \Delta^{\neq}_b + \teps_{bd} \Delta^{\neq}_d + \teps_{bc} \Delta^<_c \\
\footrel{(\ref{f4236c})}\leq
\tilde1_b \Delta^{\neq}_b + \tilde1_c \teps_{bc} \Delta^{\neq}_c
+ (\teps_{bd} + \teps_{bc}\teps_{cd}) \Delta^{\neq}_d .
\end{multline}
%
Last, bounding $\Delta^{<}_a$ requires iterating Claim~\ref{l3526} three times:
\begin{multline}\label{f4236a}
\Delta^<_a = \Delta^{[a]}_a
\leq \frac{1}{1-\epsilon_{ab}^2} \Delta^{[b]}_a + \teps_{ab} \Delta^<_b \\
\leq \frac{1}{(1-\epsilon_{ab}^2)(1-\epsilon_{ac}^2)} \Delta^{[c]}_a
+ \teps_{ac} \Delta^<_c + \teps_{ab} \Delta^<_b 
\leq \tilde1_a \Delta^{\neq a} + \teps_{ad} \Delta^<_d
+ \teps_{ac} \Delta^<_c + \teps_{ab} \Delta^<_b \\
\footrel{(\ref{f4236c},\ref{f4236b})}
\leq \tilde1_a \Delta^{\neq a} + \tilde1_b \teps_{ab} \Delta^{\neq}_b
+ \tilde1_c (\teps_{ac} + \teps_{ab}\teps_{bc}) \Delta^{\neq}_c
+ (\teps_{ad} + \teps_{ac}\teps_{cd}
+ \teps_{ab}\teps_{bd}
+ \teps_{ab}\teps_{bc}\teps_{cd}) \Delta^{\neq}_d.
\end{multline}

One can sum up Equations~(\ref{f4236d})--(\ref{f4236a}) into the matricial expression
\[\label{f5482}
\begin{pmatrix} \Delta_a^< \\ \Delta_b^< \\ \Delta_c^< \\ \Delta_d^< \end{pmatrix} \leq
\begin{pmatrix} 1 & \teps_{ab} & \teps_{ac} + \teps_{ab}\teps_{bc} &
\teps_{ad} + \teps_{ac}\teps_{cd} + \teps_{ab}\teps_{bd} + \teps_{ab}\teps_{bc}\teps_{cd} \\
0 & 1 & \teps_{bc} & \teps_{bd} + \teps_{bc}\teps_{cd} \\
0 & 0 & 1 & \teps_{cd} \\
0 & 0 & 0 & 1 \end{pmatrix}
\begin{pmatrix} \tilde1_a \Delta^{\neq}_a \\ \tilde1_b \Delta^{\neq}_b \\
\tilde1_c \Delta^{\neq}_c \\ \Delta^{\neq}_d \end{pmatrix} .\]
%
If we look back at how the square matrix in~(\ref{f5482}) has been constructed,
we find that
\[ \begin{pmatrix}
\text{\it square} \\
\text{\it matrix} \\
\text{\it in} \\
\text{\it (\ref{f5482})} \\
\end{pmatrix} =
\sum_{k=0}^{\infty} \begin{pmatrix}
0 & \teps_{ab} & \teps_{ac} & \teps_{ad} \\
0 & 0 & \teps_{bc} & \teps_{bd} \\
0 & 0 & 0 & \teps_{cd} \\
0 & 0 & 0 & 0 \end{pmatrix}^{k}
= \begin{pmatrix}
1 & -\teps_{ab} & -\teps_{ac} & -\teps_{ad} \\
0 & 1 & -\teps_{bc} & -\teps_{bd} \\
0 & 0 & 1 & -\teps_{cd} \\
0 & 0 & 0 & 1 \end{pmatrix}^{-1} ,\]
%
so in the end we obtain that
\[ \begin{pmatrix} \Delta_a^< \\ \Delta_b^< \\ \Delta_c^< \\ \Delta_d^< \end{pmatrix} \leq
M \begin{pmatrix} \Delta^{\neq}_a \\ \Delta^{\neq}_b \\ \Delta^{\neq}_c \\ \Delta^{\neq}_d
\end{pmatrix} ,\]
where $M$ is given by~(\ref{for4393}).
Then it is immediate that $\Var(f) = \sum_i \big(\Delta_i^<\big)^2 \leq
\VERT M \VERT^2 \sum_i \big(\Delta_i^{\neq}\big)^2 = \VERT M \VERT^{2} \,\mcal{E}(f,f)$, \textsc{qed}.
\end{proof}

The bound we have obtained for the spectral gap
is not symmetric by permutation of the indexes in~$I$.
It can however can be bounded by a simpler expression,
which is nearly as good as the original one in concrete situations:
%
\begin{Cor}
In Theorem~\ref{t7230}, $M$ can be replaced by the matrix
\[ M' = \begin{pmatrix} 1 & -\epsilon_{12} & \cdots & -\epsilon_{1N} \\
-\epsilon_{12} & \ddots & \ddots & \vdots \\
\vdots & \ddots & \ddots & -\epsilon_{(N-1)N} \\
-\epsilon_{1N} & \cdots & -\epsilon_{(N-1)N} & 1 \end{pmatrix}^{-1} ,\]
provided $\rho(\mbf{I}_N-M') < 1$.
\end{Cor}

\begin{proof}
Each entry of~$M$ is actually bounded by the corresponding entry of~$M'$.
To see it, we `expand' the entries of~$M$, resp.~$M'$.
First, notice that $1/(1-\epsilon_{ij}^2)$ can be expanded
into $1+\epsilon_{ij}\epsilon_{ji}
+ \epsilon_{ij}\epsilon_{ji}\epsilon_{ij}\epsilon_{ji} + \cdots$,
so that one has the expansions
\[ \tilde1_i = \sum_{i<j_1\leq\cdots\leq j_k} \prod_{\ell=1}^k
\epsilon_{ij_\ell} \epsilon_{j_\ell i} \]
and
\[ \teps_{ij} = \sum_{i<j_1\leq\cdots\leq j_k\leq j} \Big( \prod_{\ell=1}^k
\epsilon_{ij_\ell} \epsilon_{j_\ell i} \Big) \epsilon_{ij} .\]
Then, using the inversion formula $(\mrm{I}-A)^{-1} = \sum_{k=0}^\infty A^k$
for triangular arrays, one obtains that
\[ M_{ij} = \sum_{\substack{(i_0,i_1,\ldots,i_k)\\\text{\it first condition}}}
\prod_{\ell=0}^{k-1} \epsilon_{i_\ell i_{\ell+1}} ,\]
where the meaning of ``\textit{first condition}'' is given by the following
%
\begin{Def}\label{d8225}
A sequence $(i_0,\ldots,i_k)$ is said to satisfy the \textit{first condition} if:
\begin{ienumerate}
\item\label{i8226a} $i_0 = i$ and $i_k = j$;
\item\label{i8226b} $i_\ell \neq i_{\ell+1}$ for all~$\ell$;
\item $i_{\ell+1} < i_\ell$ only if $\ell \geq 1$ and $i_{\ell+1} = i_{\ell-1}$;
\item If $i_{\ell+1} < i_\ell$ and $\ell \leq k-2$,
then $i_{\ell+2} \geq i_\ell$.
\end{ienumerate}
\end{Def}

One has a similar formula for~$M'$:
\[ M'_{ij} = \sum_{\substack{(i_0,i_1,\ldots,i_k)\\\text{\it second condition}}}
\prod_{\ell=0}^{k-1} \epsilon_{i_\ell i_{\ell+1}} ,\]
where
\begin{Def}
A sequence $(i_0,\ldots,i_k)$ is said to satisfy the \textit{second condition} if
it satisfies Conditions~(\ref{i8226a}) and~(\ref{i8226b}) of Definition~\ref{d8225}.
\end{Def}

Since the \textit{second condition} is obviously weaker than the \textit{first condition},
one has $M_{ij} \leq M'_{ij}$.\linebreak[1]\strut
\end{proof}

There is a still weaker but even simpler formula:
\begin{Cor}\label{cor3632}
Defining
\[ \begin{array}{rrcl} \boldepsilon \colon & L^2(I) & \to & L^2(I) \\
& (a_i)_{i\in I} & \mapsto & \big( \sum_{j\neq i} \epsilon_{ij}a_j \big)_{i\in I} ,\end{array} \]
the spectral gap of the Glauber dynamics is at least
\[\label{fo6385} \big( 1 - \VERT \boldepsilon \VERT \big)_+^2 .\]
\end{Cor}
%
\begin{proof}
One has $M' = (\mrm{I}-\boldepsilon)^{-1}$, so, provided $\VERT A\VERT < 1$,
\[ \VERT M' \VERT = \big\VERT (\mrm{I}-A)^{-1} \big\VERT = \Big\VERT \sum_{k=0}^\infty A^k \Big\VERT
\leq \sum_{k=0}^\infty \VERT A \VERT^k = (1-\VERT A\VERT)^{-1} .\]
In the case $\VERT A\VERT \geq 1$, (\ref{fo6385}) is trivial.
\end{proof}


\subsection{Avoiding the articial phase transition}

A common situation in which we would like to apply the previous results is when $I = \ZZ^n$
and $\epsilon_{ij}$ is of the form $\epsilon(j-i)$
for some symmetric function~$\epsilon \colon \ZZ^n \to [0,1]$.
Then Corollary~\ref{cor3632} tells
that the Glauber dynamics has a (strictly) positive spectral gap as soon as
$\sum_{z\neq 0} \epsilon(z) < 1$.
But like in \S~\ref{parAvoidingThePhaseTransitionForTheToyModel},
we are going to prove that that bound is somehow `artificial'
and that it can be relaxed into the neater condition
``$\sum_{z\neq 0} \epsilon(z) < \infty$'':
%
\begin{Thm}\label{thm7431}
Suppose that $I = \ZZ^n$ and that for all~$i,j\in \ZZ^n$
one has ${\{X_i:X_j\}_\ast} \ab \leq \epsilon(j-i)$
for some symmetric function~$\epsilon \colon \ZZ^n \to [0,1]$
such that $\epsilon(z) < 1$ as soon as $z \neq 0$.
Then if $\sum_{z\in\ZZ^n} \epsilon(z) < \infty$,
the spectral gap of the Glauber dynamics is positive.
\end{Thm}

\begin{proof}
The assumption on~$\sum_z\epsilon(z)$ allows us to take $\ell < \infty$ large enough so that
\[ \sum_{z\in\ell\ZZ^n\setminus\{0\}} \epsilon(z) < 1 .\]
We split $\ZZ^n$ into a partition of~$\ell^n \eqqcolon N$ sublattices $Z_1, \ldots, Z_N$,
each lattice $Z_u$ being of the form $\ell \ZZ^n + z_u$ for some $z_u\in\ZZ^n/\ell\ZZ^n$.
Then we define an auxiliary dynamics:
%
\begin{Def}\label{def4108}
The \emph{sublattice Glauber dynamics} is the Glauber dynamics for~$\vec{X}_{\ZZ^n}$
\emph{considered as the finite-dimensional vector
$(\vec{X}_{Z_1},\ldots,\vec{X}_{Z_N})$}.
In other words, on each~$u\in\{1,\ldots,N\}$
there is an independent $\mathit{Poisson}(1)$ alarm clock,
and when clock $u$ rings, the state of the whole $\vec{X}_{Z_u}$ is flipped in one shot
according to~$\Pr\big(X_{Z_u}\big|\vec{X}_{\ZZ^n\setminus Z_u}\big)$.
\end{Def}

Now let~$f \ab {\in \ldb(\Omega)}$.
In addition to the notation of the proof of Theorem~\ref{t7230},
we introduce the following definition:
%
\begin{Def}
For~$u \in \{1,\ldots,N\}$, we define
\[ f_{(u)}^{\neq} \coloneqq f - \EE\big[ f \big| \vec{X}_{\ZZ^n \setminus Z_u} \big] .\]
\end{Def}
%
\begin{Rmk}\label{rmk8840}
The~$f_{(u)}^{\neq}$ are the equivalent of the~$f_i^{\neq}$ for the sublattice Glauber dynamics.
\end{Rmk}

Fixing some `boundary condition' $\vec{x}_{\ZZ^n\setminus Z_u}$ on~$\ZZ^n\setminus Z_u$,
we can apply Corollary~\ref{cor3632} to the Glauber dynamics
for~$\vec{X}_{Z_u}$ under the law
$\Pr\big[\Bcdot\big|\vec{X}_{\ZZ^n\setminus Z_u}=\vec{x}_{\ZZ^n\setminus Z_u}\big]$.
After integrating, one gets that
\[\label{fo4048} \Var\big( f_{(u)}^{\neq} \big) \leq
\big(1-\VERT\boldsymbol\zeta\VERT\big)^{-2} \sum_{i\in Z_u} \Var\big( f_i^{\neq} \big) ,\]
where $\boldsymbol\zeta$ is the operator on~$L^2(\ell\ZZ^n)$ defined by
\[ \big(\boldsymbol\zeta g\big)(i)
= \sum_{z\in\ell\ZZ^n\setminus\{0\}} \epsilon(z) g(i+z) ,\]
whose norm is obviously bounded by~$\sum_{z\in\ell\ZZ^n\setminus\{0\}} \eqqcolon \zeta < 1$.
Then, summing~(\ref{fo4048}) for all~$u$:
\[\label{fo8710a} \sum_{u=1}^N \Var\big( f_{(u)}^{\neq} \big) \leq
\big(1-\zeta\big)^{-2} \mcal{E}(f,f) .\]

Now, let us apply Theorem~\ref{t7230} to the sublattice Glauber dynamics [Definition~\ref{def4108}].
It yields that
\[\label{fo8710b} \Var(f) \leq \VERT M\VERT^2 \sum_{u=1}^N \Var\big( f_{(u)}^{\neq} \big) ,\]
where $M$ is some $(N\times N)$ matrix depending on the
${\{ \vec{X}_{Z_u} : \vec{X}_{Z_v} \}_*}$.
But by Theorem~\ref{lem4067},
${\{ \vec{X}_{Z_u} : \vec{X}_{Z_v} \}_*} < 1$ for all~$u\neq v$,
thus $\VERT M\VERT < \infty$ by Remark~\ref{rmk4123}.
Combining~(\ref{fo8710a}) and~(\ref{fo8710b}), we finally get that
the spectral gap of the Glauber dynamics
for~$\vec{X}_{\ZZ^n}$ is bounded below by~$\VERT M\VERT^{-2} \times (1-\zeta)^2 > 0$.\linebreak[1]\strut
\end{proof}

\begin{Rmk}
Like Theorem~\ref{lem4067}, Theorem~\ref{thm7431} could actually be stated
in the general case of `abstract' metric spaces on which some group acts profinitely.
\end{Rmk}

